\documentclass[11pt]{article}
\usepackage[a4paper,margin=1in]{geometry}
\usepackage[T1]{fontenc}
\usepackage{lmodern}
\usepackage{amsmath,amssymb,amsthm,mathtools}
\usepackage{microtype}
\usepackage{enumitem}
\usepackage{xcolor}
\usepackage[numbers,sort&compress]{natbib}
\usepackage[colorlinks=true,linkcolor=blue!55!black,citecolor=blue!55!black,urlcolor=blue!55!black]{hyperref}
\usepackage[nameinlink,noabbrev]{cleveref}

\newtheorem{theorem}{Theorem}[section]
\newtheorem{proposition}[theorem]{Proposition}
\newtheorem{lemma}[theorem]{Lemma}
\newtheorem{corollary}[theorem]{Corollary}
\theoremstyle{definition}
\newtheorem{definition}[theorem]{Definition}
\newtheorem{remark}[theorem]{Remark}
\newcommand{\E}{\mathcal E}
\newcommand{\A}{\mathfrak A}
\newcommand{\norm}[1]{\lVert#1\rVert}

\title{\textbf{Periodic Approximation of Literal Multipliers on\\
Determinant-Two M\"obius Fibers}}
\author{Liang Wang\\
\small School of Artificial Intelligence and Automation,\\[-2pt]
\small Huazhong University of Science and Technology, Wuhan 430074, P.R. China\\
\small \texttt{wangliang.f@gmail.com}}
\date{July 2026}

\begin{document}
\maketitle

\begin{abstract}
The determinant-two corridor of TPC-149 bounds every bounded periodic
multiplier on a literal two-M\"obius core, uniformly outside one
exceptional set.  We give the exact interface from that theorem to an
arbitrary literal multiplier.  The resulting bound is the sum of the
periodic-core saving and a normalized \(L^1\) approximation error.
The unpenalized error term separates over residue fibers; its
quadratic surrogate is solved by fiber means.  The supremum-norm
penalty in the full cost still couples those fibers.  This is a positive
actual-core weight-interface theorem.  It is not yet a theorem for
the missing physical occurrence registry, a generic phase or all
prefixes, and its fixed-\(X\)-power exponent is zero.
\end{abstract}

\noindent\textbf{Keywords.}
M\"obius correlation; determinant two; periodic approximation;
literal multiplier; exceptional scales.

\medskip
\noindent\fcolorbox{red!70!black}{red!3}{\parbox{0.94\linewidth}{
\textbf{Claim boundary.}
The theorem below acts on the literal determinant-two M\"obius core
once a multiplier \(w\) is supplied.  It does not construct the
frontier occurrence lift or certify that the unknown physical
multiplier has small approximation cost.  A logarithmic saving is not
a positive fixed-\(X\)-power saving and cannot pay the \(1/400\)
endpoint.  No prime-pair lower bound or twin-prime theorem is proved.}}

\section{The inherited periodic corridor}

Fix integers
\[
 a,s,R\ge1,\qquad d,u\in\mathbb Z,\qquad
 (a,s)=1,\qquad as\ {\rm odd},\qquad su-ad=2,
\]
and put
\[
 q=as,\qquad t(z)=ad+qz,\qquad
 c_z=\mu(d+sz)\mu(u+az),
\]
\[
 I_N=\{z\in\mathbb Z:N<t(z)\le2N\}.
\]
TPC-149 proves the following consequence of the quantitative affine
correlation theorem of \citet{TaoTeravainen2026}.

\begin{theorem}[TPC-149 periodic corridor]\label{thm:149}
There are absolute \(\eta_0,\kappa_0>0\) such that for every
sufficiently large \(X\) there is
\(\E_X^\star\subset[\sqrt X,X]\) satisfying
\[
 \frac1{\log X}\int_{\E_X^\star}\frac{dt}{t}
 \ll(\log X)^{-\kappa_0}.
\]
If \(qR\le(\log X)^{\eta_0}\) and
\(N\in[\sqrt X,X]\setminus\E_X^\star\), then, simultaneously for
every period-\(R\) function \(\rho\),
\begin{equation}
 \frac qN\left|\sum_{z\in I_N}c_z\rho(z)\right|
 \ll\norm{\rho}_\infty(\log X)^{-\kappa_0}.
\label{eq:periodic}
\end{equation}
\end{theorem}

The word ``simultaneously'' is essential: choosing an approximant
after seeing \(w\) does not require a union over its values.

\section{The literal-weight interface}

\begin{definition}[Periodic approximation cost]\label{def:cost}
For \(w:I_N\to\mathbb C\), define
\begin{equation}
 \A_{R,N}(w)=
 \inf_{\substack{\rho:\mathbb Z\to\mathbb C\\\rho(z+R)=\rho(z)}}
 \left\{
 \norm{\rho}_\infty(\log X)^{-\kappa_0}
 +\frac qN\sum_{z\in I_N}|w(z)-\rho(z)|
 \right\}.
\label{eq:cost}
\end{equation}
\end{definition}

\begin{theorem}[Literal-weight periodic approximation]\label{thm:main}
Under the hypotheses of \cref{thm:149},
\begin{equation}
 \boxed{\;
 \frac qN\left|\sum_{z\in I_N}c_zw(z)\right|
 \ll \A_{R,N}(w).\;}
\label{eq:main}
\end{equation}
The exceptional set is exactly the one in \cref{thm:149}.
\end{theorem}

\begin{proof}
For any admissible \(\rho\),
\[
 \sum_{I_N}c_zw(z)
 =
 \sum_{I_N}c_z\rho(z)
 +\sum_{I_N}c_z\{w(z)-\rho(z)\}.
\]
The first term is controlled by \eqref{eq:periodic}.  Since
\(|c_z|\le1\), the normalized modulus of the second is at most
\[
 \frac qN\sum_{z\in I_N}|w(z)-\rho(z)|.
\]
Take the infimum.  No exceptional-set operation occurs.
\end{proof}

\begin{corollary}[A usable promotion rule]\label{cor:promotion}
Suppose an independently source-locked physical registry supplies
\(w=w_X\), a period \(R_X\) with
\(qR_X\le(\log X)^{\eta_0}\), and an approximant satisfying
\[
 \norm{\rho_X}_\infty\ll1,\qquad
 \frac qN\sum_{I_N}|w_X-\rho_X|
 \ll(\log X)^{-\gamma}.
\]
Then the weighted actual core has a logarithmic saving with exponent
\(\min\{\kappa_0,\gamma\}\).
\end{corollary}

\begin{remark}
The corollary is an implication, not a claim that the current
frontier archive supplies \(w_X\).  That occurrence-level input
remains unavailable.
\end{remark}

\section{Residue-fiber error optimization}

For \(r\in\mathbb Z/R\mathbb Z\), write
\[
 I_{N,r}=\{z\in I_N:z\equiv r\pmod R\}.
\]

\begin{proposition}[Exact separation]\label{prop:separate}
For every \(w\),
\[
 \inf_{\rho\ {\rm period}\ R}\sum_{I_N}|w(z)-\rho(z)|
 =
 \sum_{r\bmod R}\inf_{\zeta\in\mathbb C}
 \sum_{z\in I_{N,r}}|w(z)-\zeta|.
\]
For the squared error, the unique minimizing value on every nonempty
fiber is
\[
 \bar w_r=\frac1{|I_{N,r}|}\sum_{z\in I_{N,r}}w(z).
\]
\end{proposition}

\begin{proof}
A period-\(R\) function has one independent value on each residue
fiber, proving the first identity.  Expanding
\[
 \sum_{I_{N,r}}|w(z)-\zeta|^2
 =
 \sum_{I_{N,r}}|w(z)-\bar w_r|^2
 +|I_{N,r}|\,|\zeta-\bar w_r|^2
\]
proves the second.
\end{proof}

By Cauchy--Schwarz, the fiber-mean quadratic certificate gives the
explicit upper bound
\[
 \sum_{I_N}|w-\rho|
 \le |I_N|^{1/2}
 \left(\sum_{r\bmod R}\sum_{I_{N,r}}|w-\bar w_r|^2\right)^{1/2}.
\]
Thus the interface is finite and auditable whenever literal values of
\(w\) are known.

\section{Products, phases and precise limitations}

\begin{proposition}[Exact-period products]\label{prop:product}
If \(w=\prod_{j=1}^m w_j\) and every \(w_j\) is bounded periodic of
period \(R_j\), then \(w\) is periodic of period
\(\operatorname{lcm}(R_1,\ldots,R_m)\).  It may be inserted directly
in \cref{thm:149} when
\[
 q\operatorname{lcm}(R_1,\ldots,R_m)
 \le(\log X)^{\eta_0}.
\]
\end{proposition}

Rational additive phases fall in this proposition.  A generic
additive phase need not be well approximated by a small-period
function.  TPC-158 tests that boundary exactly.  Likewise,
\cref{thm:main} treats one interval \((N,2N]\); it does not convert
exceptional-scale density into control of predetermined prefix
endpoints, the obstruction isolated in TPC-150.

\section{Audit and progress classification}

The accompanying standard-library audit source-locks the TPC-149
theorem and certificate using \texttt{CANONICAL\_UTF8\_LF\_V2},
checks exact residue-fiber quadratic minimization on rational data,
and checks the triangle decomposition on a signed fixture.  Hashes
serve only as integrity checks.

\begin{center}
\begin{tabular}{ll}
\hline
Statement & Status\\
\hline
Periodic-to-literal interface on the actual M\"obius core
  & \(\mathrm{PROVED}_{L1}\)\\
Production physical multiplier and approximation rate
  & \(\mathrm{NOT\mbox{-}TESTABLE}\)\\
Generic phase or all prefixes & not proved\\
Positive fixed-\(X\)-power saving & not proved\\
\hline
\end{tabular}
\end{center}

\section{Conclusion}

The periodic theorem now has an exact literal-weight consumer:
arithmetic cancellation is lost only through a normalized
approximation cost whose unpenalized error term is
residue-fiber-separable.  This closes a genuine
actual-core interface while exposing the next factual question:
whether the eventual occurrence registry produces small cost.  It
does not answer that question by schema, numerics or notation.

\bibliographystyle{plainnat}
\bibliography{references}
\end{document}
